% =====================================================================
% Section 2: Related Work
% =====================================================================
%
% Status: Draft §2 Related Work v6 (2026-05-18)
%   v6: advisor-driven three-bucket reorganization.
%       Bucket 1: Classic text RAG (BM25, DPR, RAPTOR, IRCoT, Self-Ask).
%       Bucket 2: GraphRAG (HippoRAG family, PropRAG, GraphRAG, HGRAG,
%                 MGRag, SentGraph, SPRiG, DFGN, Evidence Integration).
%       Bucket 3: Evidence-enhanced GraphRAG (NeocorRAG, Relink, ToG,
%                 KG-RAG).
%       Closing Summary keeps the §1 conversion-problem tie-back.
%   v5: see archive/02_related_v5_pre-advisor-rewrite_20260518.tex
%
% =====================================================================

\section{Related Work}
\label{sec:related}

\paragraph{Classic text RAG.}
Standard retrieval-augmented generation methods rank passages by
lexical or dense similarity to the question, with BM25
\citep{robertson2009bm25} and DPR \citep{karpukhin2020dpr} as the
canonical sparse and dense baselines. These methods are effective
for single-hop questions but tend to miss bridge passages whose
surface similarity to the question is low. RAPTOR
\citep{sarthi2024raptor} addresses limited cross-passage context by
recursively clustering and summarizing passages into a retrieval
tree, providing a hierarchical retrieval route over text rather
than over an explicit graph. Iterative chain-of-thought retrievers
such as IRCoT \citep{trivedi2023ircot} and Self-Ask
\citep{press2023selfask} drive multi-hop retrieval through
language-model-generated reasoning steps, alternating between
generation and lookup. These systems show that multi-hop retrieval
needs an explicit mechanism for linking evidence across documents,
but their retrieval state remains a flat or hierarchical list of
passages rather than a document-grounded evidence graph.

\paragraph{GraphRAG.}
A second line of work organizes the corpus into an explicit graph
to expose cross-document connections. HippoRAG
\citep{jimenez2024hipporag} and HippoRAG~2
\citep{gutierrez2025hipporag} extract OpenIE facts from the corpus
and run personalized PageRank over a corpus-wide graph of passage
and phrase nodes seeded from the question. PropRAG
\citep{wang2025proprag} replaces triples with propositions and
discovers multi-hop evidence through beam search over a proposition
graph, motivated by the observation that triples can lose context.
GraphRAG \citep{edge2024graphrag} builds community-summary graphs
and supports retrieval through hierarchical summary aggregation.
Hypergraph-based methods such as HGRAG \citep{hgrag2024} model
entities and passages within a hypergraph and propagate relevance
across entity-level and passage-level representations. Recent
variants explore hierarchical multi-granular graphs
\citep{hu2026mgrangrag}, sentence-level graph structures
\citep{liang2026sentgraph}, and linear CPU-only PageRank-style
retrieval \citep{wang2026sprig}. Earlier graph-based multi-hop
readers also construct document- or entity-level graphs and
integrate evidence through graph neural message passing
\citep{qiu2019dfgn,song2022evidenceintegration}. These systems
mostly differ in graph representation and search operator, ranging
from corpus-wide PageRank to beam search to hierarchical summary
aggregation. Many rely on a pre-built global structure or a summary
hierarchy, with question conditioning entering through seeds or
through a subsequent search procedure. \methodname{} builds on the
same OpenIE-based indexing line
\citep{angeli2015leveraging,gutierrez2025hipporag}, but preserves
document-grounded fact units for query-local propagation rather
than corpus-level diffusion.

\paragraph{Evidence-enhanced GraphRAG.}
A third line uses graph or KG signals to refine which evidence
reaches the reader. NeocorRAG \citep{peng2026neocorrag} keeps
documents as the reader unit but mines \emph{evidence chains} from
retrieved text through an activation-based candidate search
followed by constrained decoding, using the resulting chains to
filter and re-present passages to the reader. Relink
\citep{huang2026relink} addresses KG incompleteness by dynamically
reconstructing query-specific evidence graphs and repairing missing
facts as needed. ToG \citep{sun2024tog} performs LLM-guided
traversal over a pre-built knowledge graph, selecting reasoning
paths step by step. KAPING \citep{baek2023kaping} retrieves
subgraphs from a KG and prompts a language model with the
subgraph as context. These methods improve evidence quality by
adding chain-mining, repair, or KG-grounding on top of an
underlying retriever; they typically operate over symbolic chains
or per-question reconstructions rather than over a fixed offline
fact-source substrate.

\paragraph{Summary.}
Existing text-based RAG, GraphRAG, and evidence-enhanced GraphRAG
systems improve cross-document retrieval through richer
representations, query-time path construction, or chain-based
filtering of retrieved text. These approaches mostly target either
how to build a better corpus-level structure or how to surface
explicit reasoning chains from retrieved content. The question of
how to convert graph-level relevance into a short reader-facing
list that is both connected across documents and covering of
question-side requirements receives less direct attention.
\methodname{} addresses this conversion problem within the GraphRAG
family. It combines query-local propagation over an induced
fact-and-passage graph with relation-aware evidence admission and
coverage-aware reranking on top of an offline OpenIE graph.
